% Part: lambda-calculus
% Chapter: syntax
% Section: terms

\documentclass[../../../include/open-logic-section]{subfiles}

\begin{document}

\olfileid{lam}{syn}{trm}
\olsection{ترم‌ها}

ترم‌های حساب لامبدا به‌صورت استقرایی از ذخیره‌ای نامتناهی از متغیرهای
$\Obj{v_0}$، $\Obj{v_1}$، \dots، نماد~«$\lambd$» و پرانتزها ساخته
می‌شوند. از~$x$، $y$، $z$، \dots{} برای نامیدن متغیرها و از~$M$، $N$،
$P$، \dots{} برای نامیدن ترم‌ها استفاده می‌کنیم.

\begin{defn}[ترم‌ها] \ollabel{defn:term}
مجموعهٔ \emph{ترم‌های} حساب لامبدا به‌صورت استقرایی چنین تعریف می‌شود:
\begin{enumerate}
  \item \ollabel{defn:term-var} اگر~$x$ متغیر باشد، $x$ ترم است.
  \item \ollabel{defn:term-abs} اگر~$x$ متغیر و~$M$ ترم باشد،
    $(\lambd[x][M])$ ترم است.
  \item \ollabel{defn:term-app} اگر هر دو~$M$ و~$N$ ترم باشند،
    $(MN)$ ترم است.
\end{enumerate}
\end{defn}

اگر ترم~$(\lambd[x][M])$ مطابق~\olref{defn:term-abs} ساخته شود، می‌گوییم
حاصل یک \emph{انتزاع} است و~$x$ در~$\lambd[x]$ را \emph{!!{parameter}}
می‌نامیم. ترم~$(MN)$ که مطابق~\olref{defn:term-app} ساخته شود حاصل یک
\emph{کاربرد} است.

ترم‌هایی که در بالا تعریف شدند کاملاً پرانتزگذاری شده‌اند. این کار
می‌تواند بسیار دست‌وپاگیر شود، چنان‌که ترم
$(\lambd[x][((\lambd[x][x])(\lambd[x][(xx)]))])$ نشان می‌دهد. قراردادهایی
برای حذف پرانتزهای اضافی معرفی خواهیم کرد. بااین‌حال، تعریف رسمی تعیین
شیوهٔ ساخت یک ترم مطابق~\olref{defn:term} را آسان می‌کند. برای نمونه،
آخرین گام ساخت ترم
$(\lambd[x][((\lambd[x][x])(\lambd[x][(xx)]))])$ باید انتزاعی باشد که
در آن~!!{parameter}، $x$ است. این ترم با انتزاع از ترم
$((\lambd[x][x])(\lambd[x][(xx)]))$ به دست می‌آید، که کاربرد دو ترم است.
هر یک از این دو ترم نیز حاصل یک انتزاع است، و همین‌طور ادامه می‌یابد.

\begin{prob}
ساخت~$(\lambd[g][(\lambd[x][(g (x x))])
  (\lambd[x][(g (x x))])])$ را توصیف کنید.
\end{prob}

\end{document}
